\subsection{Constructing a method from typed intent and ordered claim rows}
\label{sec:method-specialization}

A sale method must insert the claim the caller actually names. The same
financial calculation can apply to a new customer, due date or invoice key,
yet the inserted record and the later collection pointer must change. Writing
one finished method for each saved event would merely move the native
generator's results into a catalogue. Explicit parameters instead retain the
complete method algorithm and identify exactly which values come from intent
or current state.

\code{ParameterizedCompanyMethod/1} supplies that bounded construction step.
Its template consists of literal, object, array and named-parameter nodes.
Each parameter has a declared scalar type and source: a known full-intent input,
event identifier or date, selected current-book metadata, or a pointer assembled
from literal segments and earlier text parameters. A claim key containing a
slash or tilde is escaped as a pointer segment; it is never reparsed as source.
Unknown required input remains unresolved. Duplicate, unused, forward or
unbound parameters are refused.

For a sale, parameters supply the claim key, customer, due date, event identifier
and dates in the inserted row. For a collection, they supply the selected
claim key and its remaining-balance pointer in both the precondition and the
monetary input binding. The subtraction, postings and memory write remain
visible in the full method. The input schema is compiled from its explicit
typed definition, while the state schema comes from the actual book. Complete
template, state and event identities bind the constructed request. The existing
typed frontend and method-source renderer then check the whole method.

Four templates across the two retained charts reproduce six prior sale and
collection requests and twelve new events with changed keys, customers, dates
and amounts. Comparisons cover complete requests and transition results, and
the emitted method source compiles back exactly. A deliberate edit halves the
collection algorithm's cash, receivable, flow and remaining-balance effects.
It produces exactly half the cash movement and preserves the claim balance
relation, while differing from the explicitly selected full-payment contract.
The construction mechanism therefore exposes a consequential algorithm rather
than hiding a table of fixed financial outcomes.

Assessment adds a different need: the number of writes depends on the current
claim population. \code{OrderedMethodExpansion/1} declares bounded row
expressions and an ordered list of claim keys. In the retained state-only
assessment, each row sets the assessment flag and date, copies remaining
amount into the snapshot, and records overdue amount only when the due date
is strictly earlier than the assessment date. A final write advances the
cursor. Typed expressions can read a row field, key, event metadata or known
event input, compare dates and choose between equally typed values. Both
conditional branches and every write destination are checked even when the
claim map is empty.

Order is explicit because it affects construction identity. Canonical JSON
hashing ignores the order of object fields, but the native assessment's ordered
write list follows its claim-map iteration. The new call therefore binds a
separate key-order digest and requires the list to cover every current claim
exactly once. Forward and reverse orders produce their corresponding native
requests. Their disjoint writes yield the same financial state, but that
agreement does not make the two ordered request records identical.

Eight retained and twelve new assessments match full requests and transition
results, including empty maps, multiple claims, escaped keys and due-date
boundaries. An authored inclusive comparison makes a claim overdue on its due
date. The typed transition executes that different rule, while the selected
strict lifecycle invariant rejects its snapshot. This distinction matters:
types can establish that a date comparison is meaningful, but they cannot
choose the caller's required convention.

Both constructors work with the native specializer disabled in their checks.
Their construction interface leaves full-intent eligibility separate; the next
section supplies an explicit ordered alternative for the retained contracts.
Fixed chart context, role mapping and unit-cost choices remain declared scope.
The earlier complete program composition and invariant source continue to
provide their own boundaries.

\begingroup\raggedright
Complete constructions, controls and variants are retained under
\path{data/accounting-master/parameterized-method-construction-20261026/}
and \path{data/accounting-master/ordered-assessment-expansion-20261026/}.
\par
\endgroup
